\documentclass{article}
\usepackage{amsmath,amssymb}

\begin{document}

\noindent De uitdrukking evalueert de \textit{improper integral} van een \textit{power law} functie:
\begin{equation*}
I = \int_{1}^{+\infty} \alpha x^{1-\alpha} \, dx
\end{equation*}

\noindent 1. \textit{Antiderivative} via de \textit{power rule} (geldig voor $\alpha \neq 2$):
\begin{equation*}
I = \left[ \frac{\alpha}{-\alpha + 2} x^{-\alpha + 2} \right]_{1}^{+\infty}
\end{equation*}

\noindent 2. \textit{Limit evaluation} bij het boven- en onderlimiet:
\begin{equation*}
I = \lim_{b \to +\infty} \left( \frac{\alpha}{-\alpha + 2} b^{2-\alpha} \right) - \frac{\alpha}{-\alpha + 2} (1)^{2-\alpha}
\end{equation*}

\noindent 3. \textit{Case analysis} op basis van het teken van de exponent $(2 - \alpha)$:
\begin{itemize}
    \item \textbf{Geval $\alpha > 2 \implies (2 - \alpha) < 0$:}
    \begin{equation*}
    \lim_{b \to +\infty} b^{2-\alpha} = 0 \implies I = 0 - \frac{\alpha}{-\alpha + 2} = \frac{\alpha}{\alpha - 2} \quad (\text{\textit{convergent}})
    \end{equation*}

    \item \textbf{Geval $\alpha < 2 \implies (2 - \alpha) > 0$:}
    \begin{equation*}
    \lim_{b \to +\infty} b^{2-\alpha} = +\infty \implies I = +\infty \quad (\text{\textit{divergent}})
    \end{equation*}

    \item \textbf{Geval $\alpha = 2$ (\textit{singularity} in \textit{power rule}):}
    \begin{equation*}
    I = \int_{1}^{+\infty} 2 x^{-1} \, dx = \left[ 2 \ln(x) \right]_{1}^{+\infty} = +\infty \quad (\text{\textit{divergent}})
    \end{equation*}
\end{itemize}

\noindent Dit leidt tot de piecewise functie:
\begin{equation*}
I = \begin{cases}
+\infty & \alpha \le 2 \\
\dfrac{\alpha}{\alpha - 2} & \alpha > 2
\end{cases}
\end{equation*}

\end{document}
